\chapter{مقدمه و بیان مسئله}
\label{ch:intro}

\section{مقدمه}

در بسیاری از شاخه‌های علوم پایه، مهندسی، اقتصاد، علوم کامپیوتر و مدل‌سازی پدیده‌های طبیعی، با مسائلی روبه‌رو می‌شویم که حل آن‌ها در نهایت به یافتن جواب یک معادله منجر می‌شود. بسیاری از این معادلات را می‌توان به فرم کلی زیر نوشت:
\begin{equation}
f(x)=0
\label{eq:fzero}
\end{equation}
در این رابطه، $f$ تابعی معلوم و $x$ متغیری است که باید مقدار آن به‌گونه‌ای تعیین شود که مقدار تابع برابر صفر گردد. چنین مقداری از $x$، ریشه یا جواب معادله نامیده می‌شود. اگر تابع $f$ یک تابع خطی باشد، معمولاً یافتن جواب آن ساده است و می‌توان با استفاده از روش‌های جبری مستقیم، مقدار دقیق ریشه را به دست آورد. اما در بسیاری از مسائل واقعی، تابع مورد نظر خطی نیست و ممکن است شامل عبارات چندجمله‌ای با درجه بالا، توابع مثلثاتی، نمایی، لگاریتمی یا ترکیبی از آن‌ها باشد.

در چنین حالت‌هایی، معادله حاصل یک معادله غیرخطی است. حل دقیق معادلات غیرخطی در بسیاری از موارد امکان‌پذیر نیست یا به روابط بسیار پیچیده‌ای منجر می‌شود. به همین دلیل، روش‌های عددی اهمیت زیادی پیدا می‌کنند. روش‌های عددی به جای یافتن جواب دقیق، تلاش می‌کنند با استفاده از الگوریتم‌های تکراری، دنباله‌ای از تقریب‌ها تولید کنند که به جواب واقعی مسئله نزدیک شوند.

یکی از معروف‌ترین و پرکاربردترین روش‌های عددی برای حل معادلات غیرخطی، روش نیوتن یا روش نیوتن-رافسون است. این روش با استفاده از مقدار تابع و مشتق آن در یک نقطه، تقریب بعدی ریشه را محاسبه می‌کند. ایده اصلی روش نیوتن بر پایه تقریب خطی تابع در نزدیکی نقطه فعلی است. به بیان هندسی، در هر مرحله خط مماس بر نمودار تابع رسم شده و محل برخورد این مماس با محور افقی به عنوان تقریب بعدی ریشه انتخاب می‌شود.

روش نیوتن به دلیل سرعت همگرایی بالا، سادگی فرمول و قابلیت استفاده در بسیاری از مسائل علمی و مهندسی، جایگاه ویژه‌ای در تحلیل عددی دارد. البته این روش در کنار مزایای خود، محدودیت‌هایی نیز دارد؛ از جمله وابستگی به انتخاب حدس اولیه، نیاز به وجود مشتق تابع و امکان شکست در حالتی که مشتق در نقطه‌ای صفر یا بسیار کوچک باشد. بنابراین بررسی دقیق این روش، نحوه اجرای آن، شرایط همگرایی و تحلیل نتایج حاصل از آن اهمیت زیادی دارد.

در این پروژه، روش نیوتن برای حل معادلات غیرخطی بررسی می‌شود. ابتدا مفاهیم پایه مربوط به ریشه‌یابی و معادلات غیرخطی معرفی می‌شوند. سپس فرمول روش نیوتن استخراج شده و الگوریتم آن به صورت مرحله‌به‌مرحله ارائه می‌شود. در ادامه، با استفاده از مثال‌های عددی، عملکرد روش بررسی شده و نتایج به کمک جدول‌ها و نمودارهای مناسب تحلیل می‌گردد.

\section{معادلات غیرخطی و اهمیت آن‌ها}
\label{sec:import}

معادله‌ای را غیرخطی می‌نامیم که در آن رابطه میان متغیرها به صورت خطی نباشد. به عنوان مثال، معادله زیر یک معادله خطی ساده است:
\begin{equation}
ax+b=0
\label{eq:linear}
\end{equation}
اما معادلات زیر نمونه‌هایی از معادلات غیرخطی هستند:
\begin{equation}
x^2-2=0
\label{eq:nonlinear1}
\end{equation}
\begin{equation}
\sin x - x = 0
\label{eq:nonlinear2}
\end{equation}
\begin{equation}
e^x - 3x = 0
\label{eq:nonlinear3}
\end{equation}
\begin{equation}
x^3 - x - 2 = 0
\label{eq:nonlinear4}
\end{equation}

در معادلات غیرخطی، متغیر ممکن است با توان‌های بالاتر از یک ظاهر شود یا داخل توابعی مانند سینوس، کسینوس، نمایی و لگاریتم قرار گیرد. همین ویژگی باعث می‌شود که حل چنین معادلاتی نسبت به معادلات خطی دشوارتر باشد.

اهمیت معادلات غیرخطی از آن جهت است که بسیاری از پدیده‌های واقعی ماهیت غیرخطی دارند. در واقع، خطی بودن معمولاً یک تقریب ساده‌ساز از واقعیت است. بسیاری از سیستم‌های طبیعی و مهندسی، رفتاری پیچیده و غیرخطی از خود نشان می‌دهند. برای نمونه، حرکت اجسام در میدان‌های نیرو، جریان سیالات، مدارهای الکتریکی غیرخطی، رشد جمعیت، واکنش‌های شیمیایی، مدل‌های اقتصادی و الگوریتم‌های یادگیری ماشین، همگی می‌توانند به معادلات غیرخطی منجر شوند.

در بسیاری از این مسائل، هدف یافتن مقدار یا مقادیری از متغیر است که یک رابطه تعادلی را برقرار کند. این مسئله معمولاً با حل معادله‌ای از فرم \eqref{eq:fzero} بیان می‌شود. برای مثال، در مسائل فیزیکی ممکن است مقدار تعادل یک سیستم، در مسائل اقتصادی نقطه تعادل عرضه و تقاضا، و در مسائل مهندسی مقدار بهینه یک پارامتر از طریق حل یک معادله غیرخطی به دست آید.

به همین دلیل، ریشه‌یابی معادلات غیرخطی یکی از مباحث بنیادی در ریاضیات کاربردی و تحلیل عددی محسوب می‌شود. هدف اصلی در ریشه‌یابی، یافتن مقداری مانند $\alpha$ است به‌طوری‌که:
\begin{equation}
f(\alpha)=0
\label{eq:falpha}
\end{equation}
عدد $\alpha$ ریشه تابع $f$ نامیده می‌شود. اگرچه در برخی موارد می‌توان این مقدار را به صورت دقیق محاسبه کرد، اما در بسیاری از مسائل عملی، تنها به دست آوردن یک تقریب مناسب از ریشه کافی و حتی ضروری است.

\section{ضرورت استفاده از روش‌های عددی}
\label{sec:num}

روش‌های تحلیلی و جبری، در حالت‌های ساده می‌توانند جواب دقیق یک معادله را ارائه دهند. برای مثال، معادلات درجه اول و درجه دوم معمولاً دارای فرمول‌های بسته و مشخصی برای یافتن جواب هستند. اما با افزایش پیچیدگی تابع، استفاده از روش‌های تحلیلی دشوار یا غیرممکن می‌شود.

برای نمونه، معادله زیر را در نظر بگیرید:
\begin{equation}
\cos x - x = 0
\label{eq:cosx}
\end{equation}
این معادله دارای جواب است، اما نمی‌توان ریشه آن را با روش‌های جبری ساده به دست آورد. همچنین معادلاتی مانند:
\begin{equation}
e^{-x}-x=0
\label{eq:expx}
\end{equation}
یا:
\begin{equation}
x^3+4x^2-10=0
\label{eq:cubic}
\end{equation}
ممکن است از نظر نظری قابل بررسی باشند، اما یافتن جواب دقیق آن‌ها معمولاً ساده نیست. در چنین شرایطی، روش‌های عددی ابزار اصلی برای یافتن جواب تقریبی محسوب می‌شوند.

روش‌های عددی بر اساس محاسبات تکراری عمل می‌کنند. در این روش‌ها، ابتدا یک یا چند حدس اولیه انتخاب می‌شود و سپس با استفاده از یک رابطه مشخص، تقریب‌های جدیدی برای جواب تولید می‌شود. اگر روش به درستی انتخاب شود و شرایط لازم برقرار باشد، این دنباله از تقریب‌ها به ریشه واقعی نزدیک می‌شود.

به صورت کلی، یک روش تکراری را می‌توان به شکل زیر نمایش داد:
\begin{equation}
x_{n+1}=g(x_n)
\label{eq:iteration}
\end{equation}
در این رابطه، $x_n$ تقریب فعلی و $x_{n+1}$ تقریب بعدی جواب است. تابع $g$ نیز بر اساس روش عددی مورد استفاده تعریف می‌شود.

استفاده از روش‌های عددی چند مزیت مهم دارد:
\begin{itemize}
    \item امکان حل مسائلی که جواب تحلیلی ساده ندارند.
    \item قابلیت پیاده‌سازی در رایانه و نرم‌افزارهای محاسباتی.
    \item امکان کنترل دقت جواب با تعیین معیار توقف.
    \item کاربرد گسترده در مسائل واقعی و مدل‌های علمی.
\end{itemize}

البته روش‌های عددی جواب دقیق ارائه نمی‌دهند، بلکه جواب تقریبی تولید می‌کنند. بنابراین در این روش‌ها، مفهوم خطا اهمیت زیادی دارد. معمولاً الگوریتم زمانی متوقف می‌شود که اختلاف دو تقریب متوالی کمتر از مقدار از پیش تعیین‌شده‌ای باشد. این مقدار را معمولاً با $\varepsilon$ نمایش می‌دهند و شرط توقف می‌تواند به صورت زیر نوشته شود:
\begin{equation}
|x_{n+1}-x_n|<\varepsilon
\label{eq:stopdiff}
\end{equation}
یا در برخی موارد:
\begin{equation}
|f(x_n)|<\varepsilon
\label{eq:stopfunc}
\end{equation}
در نتیجه، روش‌های عددی نه‌تنها راهکاری عملی برای حل معادلات پیچیده هستند، بلکه امکان کنترل دقت و تحلیل رفتار همگرایی را نیز فراهم می‌کنند.

\section{مروری بر روش‌های حل معادلات غیرخطی}

برای حل معادلات غیرخطی، روش‌های عددی متعددی وجود دارد. هر یک از این روش‌ها دارای مزایا، محدودیت‌ها و شرایط کاربرد خاص خود هستند. انتخاب روش مناسب معمولاً به نوع تابع، اطلاعات موجود از تابع، دقت مورد نیاز و هزینه محاسباتی بستگی دارد.

یکی از ساده‌ترین روش‌ها، روش دوبخشی است. در این روش ابتدا بازه‌ای انتخاب می‌شود که تابع در دو سر آن علامت‌های متفاوت داشته باشد. سپس بازه به دو قسمت تقسیم شده و بخشی که شامل ریشه است انتخاب می‌شود. این فرایند تا رسیدن به دقت مطلوب ادامه پیدا می‌کند. روش دوبخشی همگرایی تضمین‌شده‌ای دارد، اما سرعت آن نسبتاً کم است.

روش دیگر، روش نابجایی\footnote{\lr{regula falsi}} است. این روش نیز مانند روش دوبخشی از تغییر علامت تابع در یک بازه استفاده می‌کند، اما به جای نصف کردن بازه، از خط واصل دو نقطه انتهایی برای تخمین محل ریشه استفاده می‌کند. این روش معمولاً از روش دوبخشی سریع‌تر است، اما در برخی موارد ممکن است سرعت همگرایی آن کاهش یابد.

روش سکانت یکی دیگر از روش‌های پرکاربرد است. این روش شبیه روش نیوتن است، اما به جای استفاده از مشتق دقیق تابع، مشتق را با استفاده از دو نقطه متوالی تقریب می‌زند. مزیت روش سکانت این است که نیازی به محاسبه مشتق ندارد، اما ممکن است نسبت به روش نیوتن از نظر تعداد تکرارها کندتر باشد.

در میان این روش‌ها، روش نیوتن یکی از سریع‌ترین و شناخته‌شده‌ترین روش‌ها برای ریشه‌یابی است. این روش از مشتق تابع برای ساخت یک تقریب خطی استفاده می‌کند و در شرایط مناسب، دارای همگرایی درجه دوم است. همگرایی درجه دوم به این معناست که اگر تقریب فعلی به اندازه کافی به ریشه واقعی نزدیک باشد، تعداد ارقام صحیح جواب در هر مرحله تقریباً دو برابر می‌شود.

با وجود این مزیت مهم، روش نیوتن نیازمند محاسبه مشتق تابع است و همچنین انتخاب حدس اولیه در موفقیت آن نقش مهمی دارد. اگر حدس اولیه مناسب نباشد، ممکن است روش به ریشه مورد نظر همگرا نشود یا حتی واگرا گردد. به همین دلیل، تحلیل روش نیوتن و بررسی شرایط اجرای آن اهمیت زیادی دارد.

\section{معرفی کلی روش نیوتن}

روش نیوتن که گاهی با نام روش نیوتن-رافسون نیز شناخته می‌شود، یک روش تکراری برای حل معادلات غیرخطی از فرم \eqref{eq:fzero} است. ایده اصلی این روش بر پایه استفاده از خط مماس بر نمودار تابع در نقطه فعلی است.

فرض کنید $x_n$ تقریب فعلی ریشه باشد. اگر تابع $f$ در نزدیکی این نقطه مشتق‌پذیر باشد، می‌توان تابع را در اطراف $x_n$ با یک خط مماس تقریب زد. معادله خط مماس بر نمودار تابع در نقطه $(x_n,f(x_n))$ به صورت زیر نوشته می‌شود:
\begin{equation}
y=f(x_n)+f'(x_n)(x-x_n)
\label{eq:tangent}
\end{equation}

در روش نیوتن، نقطه برخورد این خط مماس با محور $x$ به عنوان تقریب بعدی ریشه انتخاب می‌شود. برای یافتن این نقطه، کافی است مقدار $y$ را برابر صفر قرار دهیم:
\begin{equation}
0=f(x_n)+f'(x_n)(x_{n+1}-x_n)
\label{eq:tangentzero}
\end{equation}
با حل این رابطه نسبت به $x_{n+1}$، فرمول اصلی روش نیوتن به دست می‌آید:
\begin{equation}
x_{n+1}=x_n-\frac{f(x_n)}{f'(x_n)}
\label{eq:newton}
\end{equation}

این رابطه نشان می‌دهد که برای محاسبه تقریب بعدی، باید مقدار تابع و مشتق آن در نقطه فعلی محاسبه شود. سپس مقدار جدید به عنوان تقریب بهتر ریشه در مرحله بعد استفاده می‌شود.

از دیدگاه هندسی، اگر $x_n$ به ریشه واقعی نزدیک باشد و تابع رفتار مناسبی داشته باشد، خط مماس معمولاً محور $x$ را در نقطه‌ای نزدیک‌تر به ریشه قطع می‌کند. بنابراین با تکرار این فرایند، دنباله‌ای از تقریب‌ها به دست می‌آید:
\[
x_0, x_1, x_2, \dots, x_n
\]
که انتظار می‌رود به ریشه واقعی $\alpha$ همگرا شود.

روش نیوتن در بسیاری از مسائل عملکرد بسیار سریعی دارد. اما این سرعت بالا زمانی حاصل می‌شود که شرایطی مانند مشتق‌پذیری تابع، غیرصفر بودن مشتق در نزدیکی ریشه و انتخاب حدس اولیه مناسب برقرار باشد. در غیر این صورت، روش ممکن است با مشکل مواجه شود.

\section{بیان مسئله}
\label{sec:state}

مسئله اصلی این پروژه، بررسی روش نیوتن برای یافتن جواب تقریبی معادلات غیرخطی است. به طور مشخص، فرض می‌شود تابعی مانند $f$ داده شده است و هدف یافتن عددی مانند $\alpha$ است، به‌گونه‌ای که رابطه \eqref{eq:falpha} برقرار باشد.

در بسیاری از موارد، مقدار دقیق $\alpha$ مشخص نیست یا به‌دست‌آوردن آن با روش‌های تحلیلی ساده امکان‌پذیر نیست. بنابراین لازم است از یک روش عددی استفاده شود تا با شروع از یک حدس اولیه $x_0$، دنباله‌ای از تقریب‌ها تولید شود که به ریشه واقعی نزدیک شوند.

در روش نیوتن، این دنباله با رابطه \eqref{eq:newton} ساخته می‌شود.

مسئله این است که چگونه می‌توان با استفاده از این رابطه، ریشه معادله غیرخطی را با دقت مطلوب محاسبه کرد و رفتار الگوریتم را از نظر همگرایی، خطا و تعداد تکرارها بررسی نمود.

در این پروژه تلاش می‌شود به پرسش‌های زیر پاسخ داده شود:
\begin{itemize}
    \item روش نیوتن چگونه از نظر هندسی و تحلیلی به دست می‌آید؟
    \item مراحل اجرای الگوریتم نیوتن برای حل معادلات غیرخطی چیست؟
    \item معیار توقف مناسب برای این روش چگونه تعریف می‌شود؟
    \item انتخاب حدس اولیه چه تأثیری بر عملکرد روش دارد؟
    \item روش نیوتن در یک مثال عددی مشخص با چه سرعتی به جواب نزدیک می‌شود؟
    \item چه محدودیت‌هایی ممکن است باعث شکست یا کاهش کارایی روش نیوتن شوند؟
\end{itemize}

برای پاسخ به این پرسش‌ها، ابتدا مبانی نظری روش بررسی می‌شود. سپس الگوریتم نیوتن به صورت شبه‌کد ارائه شده و مراحل آن با استفاده از فلوچارت نمایش داده می‌شود. در ادامه، روش روی یک یا چند تابع نمونه اجرا شده و نتایج حاصل در قالب جدول و نمودار تحلیل می‌گردد.

\section{اهداف پروژه}

هدف کلی این پروژه، بررسی و پیاده‌سازی روش نیوتن برای حل عددی معادلات غیرخطی است. برای رسیدن به این هدف کلی، چند هدف جزئی در نظر گرفته شده است.

نخستین هدف، معرفی مفهوم ریشه‌یابی و اهمیت حل معادلات غیرخطی در مسائل علمی و مهندسی است. بدون شناخت کافی از مسئله ریشه‌یابی، درک اهمیت روش‌های عددی دشوار خواهد بود.

هدف دوم، استخراج فرمول روش نیوتن و توضیح ایده هندسی آن است. در این بخش نشان داده می‌شود که چگونه استفاده از خط مماس بر نمودار تابع می‌تواند به یک رابطه تکراری برای محاسبه ریشه منجر شود.

هدف سوم، ارائه الگوریتم روش نیوتن به صورت مرحله‌به‌مرحله است. این الگوریتم شامل انتخاب حدس اولیه، محاسبه مقدار تابع و مشتق، به‌روزرسانی تقریب ریشه و بررسی شرط توقف است.

هدف چهارم، بررسی عملکرد روش با استفاده از مثال‌های عددی است. در این مرحله، الگوریتم روی تابعی مشخص اجرا می‌شود و مقادیر به‌دست‌آمده در هر تکرار در قالب جدول نمایش داده می‌شوند.

در نهایت هدف، تحلیل رفتار همگرایی روش نیوتن است. در این بخش کاهش خطا، تعداد تکرارها و تأثیر حدس اولیه مورد بررسی قرار می‌گیرد.

\section{ساختار کلی پروژه}

این پروژه در چند فصل تنظیم شده است تا موضوع به صورت تدریجی و منظم بررسی شود.

در فصل اول، مقدمه‌ای درباره اهمیت حل معادلات غیرخطی و ضرورت استفاده از روش‌های عددی ارائه می‌شود. همچنین مسئله اصلی پروژه، اهداف و ساختار کلی گزارش بیان می‌گردد.

در فصل دوم، مبانی نظری روش نیوتن بررسی می‌شود. در این فصل ابتدا مفهوم ریشه تابع و تقریب عددی معرفی شده و سپس فرمول روش نیوتن از دیدگاه هندسی و تحلیلی استخراج می‌شود. همچنین شرایط همگرایی و برخی محدودیت‌های نظری روش توضیح داده می‌شود.

در فصل سوم، الگوریتم نیوتن به صورت دقیق و مرحله‌به‌مرحله ارائه می‌شود. در این فصل، شبه‌کد الگوریتم، معیار توقف، فلوچارت اجرای روش و نکات مربوط به پیاده‌سازی بررسی می‌شوند.

در فصل چهارم، روش نیوتن روی مثال‌های عددی اجرا می‌شود. نتایج حاصل در قالب جدول‌ها و نمودارهای مختلف نمایش داده شده و رفتار همگرایی روش تحلیل می‌گردد. همچنین تأثیر حدس اولیه و مقدار خطا بر عملکرد روش بررسی می‌شود.
